\documentclass[12pt]{article}
\usepackage{amsmath,amssymb,amsfonts}
\usepackage[margin=1in]{geometry}

\begin{document}

\textbf{Chapter 6, Problem 39.} Show that:
\begin{itemize}
\item[(a)] $\Gamma(x+1) = x\Gamma(x)$ for $x > 0$
\item[(b)] $\Gamma(n+1) = n!$, where $n$ is a positive integer.
\end{itemize}

\bigskip
\textbf{Solution.}

The Gamma function is defined by:
\[
\Gamma(x) = \int_0^{\infty}t^{x-1}e^{-t}\,dt, \qquad x > 0.
\]

\textbf{(a)} Integrate by parts with $u = t^x$ and $dv = e^{-t}\,dt$:
\[
\Gamma(x+1) = \int_0^{\infty}t^{x}e^{-t}\,dt = \left[-t^x e^{-t}\right]_0^{\infty} + \int_0^{\infty}xt^{x-1}e^{-t}\,dt.
\]

The boundary term: as $t \to \infty$, $t^x e^{-t} \to 0$ (exponential beats polynomial). At $t = 0$: $0^x \cdot 1 = 0$ (for $x > 0$). So the boundary term vanishes.

\[
\Gamma(x+1) = x\int_0^{\infty}t^{x-1}e^{-t}\,dt = x\,\Gamma(x). \quad
\]

\textbf{(b)} From part (a), applying the recursion repeatedly:
\[
\Gamma(n+1) = n\,\Gamma(n) = n(n-1)\Gamma(n-1) = \cdots = n(n-1)(n-2)\cdots 2 \cdot 1 \cdot \Gamma(1).
\]

Evaluate $\Gamma(1)$:
\[
\Gamma(1) = \int_0^{\infty}t^{0}e^{-t}\,dt = \int_0^{\infty}e^{-t}\,dt = 1.
\]

Therefore:
\[
\Gamma(n+1) = n! \cdot 1 = n!. \quad
\]

\end{document}
